\Needspace{6\baselineskip}
\begin{center}\hypertarget{native-254-article181}{}\pdfbookmark[2]{181.}{article-181}181.\end{center}
Quoties itaque unica tantum incognita $p$ adest, ad cuius determinationem
valores functionum $ap+n$, $a'p+n'$, $a''p+n''$ etc. resp. inventi sunt $=M$,
$M'$, $M''$ etc. et quidem per observationes aeque exactas, valor maxime
probabilis ipsius $p$ erit
\[
=\frac{am+a'm'+a''m''+\text{etc.}}{aa+a'a'+a''a''+\text{etc.}}=A,
\]
scribendo $m$, $m'$, $m''$ etc. resp. pro $M-n$, $M'-n'$, $M''-n''$ etc.

Iam ut gradus praecisionis in hoc valore praesumendae aestimetur,
supponemus, probabilitatem erroris $\Delta$ in observationibus exprimi per
$\frac{h}{\BareRoot{\pi}}e^{-hh\Delta\Delta}$. Hinc probabilitas, valorem
verum ipsius $p$ esse $=A+p'$, proportionalis erit functioni
\[
e^{-hh((ap-m)^2+(a'p-m')^2+(a''p-m''^2+\text{etc.})},
\]
si pro $p$ substituitur $A+p'$. Exponens huius functionis reduci potest ad
formam $-hh(aa+a'a'+a''a''+\text{etc.})(pp-2pA+B)$, ubi $B$ a $p$
independens est: proin functio ipsa proportionalis erit huic
\[
e^{-hh(aa+a'a'+a''a''+\text{etc.})p'p'}.
\]
Patet itaque, valori $A$ eundem praecisionis gradum tribuendum esse, ac si
inventus esset per observationem immediatam, cuius praecisio ad
praecisionem observationum primitivarum esset ut
$h\BareRoot{(aa+a'a'+a''a''+\text{etc.})}$ ad $h$, sive ut
$\BareRoot{(aa+a'a'+a''a''+\text{etc.})}$ ad $1$.
